Para células geometricamente semelhantes, a razão superfície/massa varia como o inverso do comprimento característico $L$:
\[
\frac{S}{m} \propto \frac{L^2}{L^3} = \frac{1}{L} \propto V^{-1/3}.
\]
Para que $S/m$ permaneça \textbf{constante}, a célula não pode crescer antes de dividir.

\textbf{Cálculo para o pepino (divisão a $1{,}5\times$ o volume):}
Se a célula cresce de $V_0$ para $V_1 = 1{,}5\,V_0$:
\[
L_1 = L_0\,1{,}5^{1/3} \approx 1{,}145\,L_0.
\]
A superfície e a massa escaladas:
\[
S_1 = S_0\,1{,}5^{2/3} \approx 1{,}31\,S_0, \qquad m_1 = 1{,}5\,m_0.
\]
Razão superfície/massa no momento da divisão:
\[
\frac{S_1}{m_1} = \frac{1{,}31\,S_0}{1{,}5\,m_0} = 0{,}875\,\frac{S_0}{m_0}.
\]
A razão $S/m$ \textbf{diminuiu} cerca de $12{,}5\%$ em relação ao valor inicial, em vez de permanecer constante.

\textbf{Conclusão:} esse comportamento não corresponde ao padrão descrito como normal. Um pepino que divide suas células quando elas atingem $1{,}5\times$ o volume das demais não mantém $S/m$ constante, portanto \textbf{não é um pepino normal} segundo esse critério.

